
\begin{answerBox}{2in}          % please don't edit this line, doing so will throw off the page layout

  %%%%%%%%%%%%%%%%%%%%%%%%%%%%%%%%%%%%%%%%%%%%%%%%%%%%%%
  %%%%%%%%%%%%%%%%%%%%%%%%%%%%%%%%%%%%%%%%%%%%%%%%%%%%%%
  %%
  %% PROBLEM 1 COLLABORATORS AND CITATIONS 
  %%
  %%%%%%%%%%%%%%%%%%%%%%%%%%%%%%%%%%%%%%%%%%%%%%%%%%%%%%
  %%%%%%%%%%%%%%%%%%%%%%%%%%%%%%%%%%%%%%%%%%%%%%%%%%%%%%
  
  \paragraph{Problem \arabic{problem} collaboration and citations}~\GradescopeComment
  {This should be on page 2 (bottom) of your PDF.}

  I collaborated on this problem with the following people:
  \begin{blue}
    \begin{itemize}
    \item \REPLACEME            % write "none"if none 
    \end{itemize}
  \end{blue}

  My answers use ideas from the following sources (other than the lecture notes and textbooks): 
  \begin{blue}
    \begin{itemize}
    \item \REPLACEME            % write "none"if none 
    \end{itemize}
  \end{blue}

\end{answerBox}

\newpage

\paragraph{Problem 1 answer}~\GradescopeComment{This should be on page 3 (top) of your PDF.}

\noindent\begin{answerBox}{1.75in}        % please don't edit this line, doing so will throw off the page layout

  %%%%%%%%%%%%%%%%%%%%%%%%%%%%%%%%%%%%%%%%%%%%%%%%%%%%%% 
  %%%%%%%%%%%%%%%%%%%%%%%%%%%%%%%%%%%%%%%%%%%%%%%%%%%%%% 
  %% 
  %% PROBLEM 1(a) ANSWER 
  %% 
  %%%%%%%%%%%%%%%%%%%%%%%%%%%%%%%%%%%%%%%%%%%%%%%%%%%%%% 
  %%%%%%%%%%%%%%%%%%%%%%%%%%%%%%%%%%%%%%%%%%%%%%%%%%%%%% 

  \paragraph{(a)}~{}
  \begin{lemma}
    For any $i\in\{0,1,\ldots,n\}$ and $t\in\{0,1,\ldots, T\}$, for $N$ as defined in the problem statement,

    \[
      N(i, t) =
      \begin{blue}
        \begin{cases}
          \REPLACEME & \text{ if } \REPLACEME, \\
          \REPLACEME & \text{ if } \REPLACEME, \\
          \REPLACEME & \text{ if } \REPLACEME, \\
          \REPLACEME & \text{ if } \REPLACEME. \\
        \end{cases}
      \end{blue}
    \]

    %% For reference, here is the LaTeX source for the recurrence relation for Knapsack from Section 7.6:
    %
    % \[
    %   V(i,\ell) =
    %   \begin{cases}
    %     0 & \text{if } i=0,
    %     \\
    %     \max(  V(i-1, \ell), V(i-1, \ell-w_i) + v_i) & \text{if } i > 0 \text{ and } w_i \le \ell,
    %     \\
    %     V(i-1, \ell) & \text{if } i > 0 \text{ and } w_i > \ell.
    %   \end{cases}
    % \]

  \end{lemma}
\end{answerBox}

\begin{answerBox}{7in}   % please don't edit this line, doing so will throw off the page layout
  
  %%%%%%%%%%%%%%%%%%%%%%%%%%%%%%%%%%%%%%%%%%%%%%%%%%%%%% 
  %%%%%%%%%%%%%%%%%%%%%%%%%%%%%%%%%%%%%%%%%%%%%%%%%%%%%% 
  %% 
  %% PROBLEM 1(b) ANSWER 
  %% 
  %%%%%%%%%%%%%%%%%%%%%%%%%%%%%%%%%%%%%%%%%%%%%%%%%%%%%% 
  %%%%%%%%%%%%%%%%%%%%%%%%%%%%%%%%%%%%%%%%%%%%%%%%%%%%%% 

  \paragraph{(b)}

  \begin{longFormProof}%[First long-form proof.]

    \step Consider any $i\in\{0,\ldots,n\}$ and $t\in\{0,1,\ldots,T\}$.

    % [review 2026-09-27] fix: J: summation index i shadowed the outer i of N(i,t); renamed to j
    \step Recall that $N(i, t) = \big|\big\{ S\subseteq \{1,2,\ldots,i\} : \sum_{j\in S} v_{j} = t\big\}\big|$.

    \medskip 

    \begin{block}
      \step \underline{Case 1.} Consider the first base case \BLUE{\REPLACEME}.

      \begin{blue}
        
        \step \REPLACEME

      \end{blue}

      \step So $N(i, t) = \BLUE{\REPLACEME}$.
   
    \end{block}

    \medskip

    \begin{block}
      \step \underline{Case 2.} Consider the second base case \BLUE{\REPLACEME}.

      \begin{blue}
        
        \step \REPLACEME

      \end{blue}

      \step So $N(i, t) = \BLUE{\REPLACEME}$.
   
    \end{block}

    \medskip

    \begin{block}
      \step \underline{Case 3.}  Next consider the case when \BLUE{\REPLACEME}.

      \step Consider the subsets $S\subseteq \{1,2,\ldots,i\}$ such that $v(S) = t$.

      \step Partition these into two types: (A) those that don't include $i$, and (B) those that do.

      \step\label{step: A or B} $N(i, t)$ is the number of subsets of either type.

      \smallskip 
      
      \step The type-A subsets are just the subsets of $\{1,\ldots,i-1\}$ with $v(S) = t$.

      \begin{blue}
        
        \step\label{step: A}
        \REPLACEME

        \smallskip

        \step The type-B subsets \REPLACEME

        \step\label{step: B}  \REPLACEME
        \smallskip 

        \step By Steps~\ref{step: A or B},~\ref{step: A} and~\ref{step: B}, \REPLACEME

      \end{blue}

      \step So the recurrence holds in this case.
    \end{block}

    \medskip 

    \begin{block}
      % [review 2026-09-27] fix: N: Case 4 placeholder was not wrapped in \BLUE like Case 3's
      \step \underline{Case 4.}  Finally consider the remaining case, when \BLUE{\REPLACEME}.

      \begin{blue}
        
        \step \REPLACEME

        \step \REPLACEME

        \step \REPLACEME

      \end{blue}

      \step So the recurrence holds in this case. 

    \end{block}

    \medskip 

    \step By the case analysis (Cases 1--4) above, the recurrence holds.

  \end{longFormProof}

\end{answerBox}

%%  THE COLLABORATORS AND CITATIONS STATEMENT IS AT THE TOP OF THIS FILE


%%% Local Variables:
%%% mode: latex
%%% TeX-master: "homework_9"
%%% End:
